\documentclass[aspectratio=169,10pt]{beamer}
\usepackage[utf8]{inputenc}
\usepackage[T1]{fontenc}
\usepackage[brazil]{babel}
\usepackage{amsmath,amssymb}
\usepackage{booktabs}
\usepackage{array}
\usepackage{graphicx}
\usepackage{xcolor}
\usepackage{tikz}
\usepackage{siunitx}
\usetheme{Madrid}
\usecolortheme{default}
\setbeamertemplate{navigation symbols}{}
\setbeamertemplate{caption}[numbered]
\setbeamerfont{caption}{size=\scriptsize}
\setbeamerfont{frametitle}{size=\large}
\graphicspath{{../}}

\title[Guia de fala]{Guia de fala --- Disserta\c{c}\~ao de Mestrado}
\subtitle{LLMs via Otimiza\c{c}\~ao Inversa --- Execu\c{c}\~ao produ\c{c}\~ao 30/06--01/07/2026}
\author{George Gomes da Silva}
\date{Julho de 2026}

\begin{document}
\shorthandoff{"}

\begin{frame}{Slide 1 --- Capa}
\footnotesize
Bom dia, professor. Vou apresentar minha disserta\c{c}\~ao de mestrado --- \emph{Explicando Decis\~oes de LLMs via Otimiza\c{c}\~ao Inversa} --- orientada pelo Prof.\ Raul Fonseca Neto. A apresenta\c{c}\~ao est\'a organizada em tr\^es blocos auto-contidos: no Bloco 1 o LLM \'e a fonte de decis\~oes; no Bloco 2 o LLM \'e o aprendiz (Problema E); no Bloco 3 aplicamos tudo na base real peso × altura. Os n\'umeros desta apresenta\c{c}\~ao v\^em de execu\c{c}\~ao completa de produ\c{c}\~ao com 3 sementes e 3 repeti\c{c}\~oes.
\end{frame}

\begin{frame}{Slide 2 --- Motiva\c{c}\~ao}
\footnotesize
LLMs s\~ao caixas-pretas. Pedir explica\c{c}\~ao em linguagem natural d\'a racionaliza\c{c}\~ao, n\~ao causa real. Em aplica\c{c}\~oes cr\'iticas isso \'e insuficiente. A pergunta central: o LLM tem um processo decis\'orio est\'avel e aprend\'ivel? Em vez de pedir explica\c{c}\~ao, observamos as decis\~oes e recuperamos o crit\'erio por tr\'as delas via otimiza\c{c}\~ao inversa.
\end{frame}

\begin{frame}{Slide 3 --- Hip\'oteses}
\footnotesize
Cinco hip\'oteses. H1 e H2 no Bloco 1: crit\'erio est\'avel e few-shot amplifica. H3 e H4 no Bloco 2: LLM aprende via in-context learning e exemplos hard deveriam ensinar mais que easy. H5 transversal. Os achados desta execu\c{c}\~ao mostram que H4 \'e refutada em n\_shot pequeno mas revertida em n\_shot=40 --- diagn\'ostico mais nuan\c{c}ado que o esperado.
\end{frame}

\begin{frame}{Slide 4 --- Abordagem}
\footnotesize
Trato cada classifica\c{c}\~ao do LLM como sa\'ida de um classificador nearest-centroid ponderado por Mahalanobis diagonal. Dados pares ponto-r\'otulo do LLM, encontro $W \geq 0$ que reproduz suas decis\~oes. Uso dois algoritmos independentes --- Perceptron Estruturado e NNLS. Se ambos chegam ao mesmo $W$, a evid\^encia n\~ao depende do m\'etodo.
\end{frame}

\begin{frame}{Slide 5 --- Defini\c{c}\~ao matem\'atica}
\footnotesize
A dist\^ancia \'e $d_W(x,c) = \sum w_j (x_j - c_j)^2$. Classificador: argmin sobre as classes. Mahalanobis diagonal corresponde a $\Sigma^{-1}$ quando features s\~ao independentes. $O(d)$ par\^ametros e interpretabilidade direta. Destaco a nomenclatura das fases --- Problema D \'e meia-lua n\~ao-linear (Bloco 1), Problema E \'e LLM aprendiz (Bloco 2).
\end{frame}

\begin{frame}{Slide 6 --- Configura\c{c}\~oes}
\footnotesize
Esta \'e a execu\c{c}\~ao de produ\c{c}\~ao completa: 3 sementes (42, 123, 7), 3 repeti\c{c}\~oes por configura\c{c}\~ao, n\_shots de 0 a 40, 4 estrat\'egias, 3 experts. 9 observa\c{c}\~oes por c\'elula para bootstrap CI. Modelo gpt-4o-mini com temperatura zero. Foram 119.713 chamadas ao LLM ao todo, com taxa de fallback do parser de 0,00\% --- ou seja, nenhum r\'otulo foi decidido por hash, o que sustenta a credibilidade dos kappas. Rodamos tamb\'em uma auditoria offline sobre todas as intera\c{c}\~oes registradas, e ela mediu o n\~ao-determinismo da temperatura zero: 4,753\% das 29.116 consultas repetidas --- 1.384 delas --- receberam respostas divergentes entre repeti\c{c}\~oes. Esse n\'umero \'e importante porque \'e o piso de ru\'ido da nossa medi\c{c}\~ao: nenhuma m\'etrica de concord\^ancia com o LLM pode passar de uns 95\%, e \'e contra esse teto que os kappas e as fidelidades devem ser lidos. A coleta ocorreu entre 30 de junho e primeiro de julho de 2026, pelo alias gpt-4o-mini --- sem snapshot datado do provedor, o que registramos como limita\c{c}\~ao. Hiperpar\^ametros do Perceptron: $\eta=0{,}001$ e $C=1{,}0$, conforme CILAMCE 2017.
\end{frame}

\begin{frame}{Slide 7 --- Algoritmo 1: Perceptron Estruturado}
\footnotesize
Mostro o pseudoc\'odigo completo do Perceptron Estruturado, do artigo de Coelho, Borges e Fonseca Neto, CILAMCE 2017, Se\c{c}\~ao 5.2. A estrutura tem duas camadas: busca bin\'aria no $\gamma$ \'otimo (loop externo, Equa\c{c}\~ao 31) e Perceptron interno que corrige $w$ por ponto violado (Equa\c{c}\~ao 29). Os hiperpar\^ametros do artigo s\~ao $\eta=0{,}001$ e $C$ entre 0{,}1 e 1{,}0; usamos 1{,}0. Pe\c{c}o aten\c{c}\~ao para a proje\c{c}\~ao em $w \geq 0$ que garante semidefini\c{c}\~ao positiva da matriz diagonal. O atributo \texttt{gamma\_history} capturado em cada itera\c{c}\~ao fornece o diagn\'ostico da busca bin\'aria que o professor pediu na reuni\~ao de 30/04 em torno de $\sim$520s --- a figura aparece no slide 83.
\end{frame}

\begin{frame}{Slide 8 --- Algoritmo 2: NNLS}
\footnotesize
Mostro o pseudoc\'odigo do NNLS de Lawson e Hanson de 1974. Constru\c{c}\~ao da matriz $A$ com diferen\c{c}as de quadrados de dist\^ancias aos centr\'oides e vetor $b$ de uns (margem-alvo). Chamada \'unica para \texttt{scipy.optimize.nnls(A, b)} resolve a otimiza\c{c}\~ao $\min \|Aw-b\|^2$ com $w \geq 0$. Vantagens sobre Perceptron: zero hiperpar\^ametros, determin\'istico (reproduz\'ivel bit-a-bit), otimalidade global garantida do problema convexo QP. Cosseno t\'ipico com o $W$ do Perceptron acima de 0{,}97 --- valida\c{c}\~ao cruzada do m\'etodo.
\end{frame}

\begin{frame}{Slide 9 --- Robustez da coleta}
\footnotesize
Pipeline com 5 camadas de prote\c{c}\~ao: parser de 7 camadas, at\'e 5 retentativas para respostas malformadas, fallback determin\'istico por hash MD5 sem vi\'es geom\'etrico, concorr\^encia ass\'incrona com sem\'aforo de 10 que reduz tempo em uma ordem de magnitude, timeout de 60s com wrap de 90s contra travamento.
\end{frame}

\begin{frame}{Slide 10 --- Bloco 1 --- Transi\c{c}\~ao Bloco 1}
\footnotesize
Entramos no Bloco 1, onde o LLM \'e a fonte de decis\~oes. O LLM rotula, n\'os aprendemos a m\'etrica $W$. Cobrir\'e os problemas A, B, C lineares e o D meia-lua. Hip\'oteses: H1 e H2.
\end{frame}

\begin{frame}{Slide 11 --- Bloco 1 --- Vis\~ao geral dos 4 problemas (apenas texto)}
\footnotesize
Apresento textualmente os 4 problemas que ser\~ao detalhados nos pr\'oximos slides. Problema A \'e a bancada de treinamento, gaussianas lineares com centr\'oides em $(-2, 0)$ e $(2, 0)$, $n=150$. Problema B \'e rota\c{c}\~ao hor\'aria agressiva, centr\'oides em $(-2; +1{,}5)$ e $(2; -1{,}5)$, $n=100$. Problema C \'e a rota\c{c}\~ao anti-hor\'aria pareada simetricamente, centr\'oides em $(-2; -1{,}5)$ e $(2; +1{,}5)$. Problema D \'e a meia-lua sint\'etica do \texttt{make\_moons} com $n=150$ e ru\'ido 0{,}15 --- problema-controle para R3/R4. Os primeiros 3 s\~ao lineares com transfer\^encia desafiadora; o 4\textordmasculine{} \'e o teste expl\'icito de n\~ao-linearidade.
\end{frame}

\begin{frame}{Slide 12 --- Bloco 1 --- Problema A}
\footnotesize
A bancada de treinamento. \texttt{make\_blobs}, centr\'oides em $(-2, 0)$ e $(2, 0)$, desvio padr\~ao 1{,}2 isotr\'opico, 150 amostras. Linearmente separ\'avel em $x_1$. F\'acil para o LLM.
\end{frame}

\begin{frame}{Slide 13 --- Bloco 1 --- Problema B}
\footnotesize
Testa transfer\^encia em geometria deslocada agressivamente. Centr\'oides em $(-2; +1{,}5)$ e $(2; -1{,}5)$ --- rota\c{c}\~ao hor\'aria. Magnitude $\pm 1{,}5$, simetricamente pareada com C. Mais agressiva que a vers\~ao anterior de $\pm 0{,}8$.
\end{frame}

\begin{frame}{Slide 14 --- Bloco 1 --- Problema C}
\footnotesize
Mesma magnitude de B mas dire\c{c}\~ao oposta: anti-hor\'aria. Centr\'oides em $(-2; -1{,}5)$ e $(2; +1{,}5)$. O orientador pediu esta orienta\c{c}\~ao oposta na reuni\~ao de 30 de abril, porque o C anterior parecia c\'opia de B. Agora B e C s\~ao pares: mesma severidade, dire\c{c}\~oes opostas.
\end{frame}

\begin{frame}{Slide 15 --- Bloco 1 --- Problema D meia-lua}
\footnotesize
Meia-lua entra no Bloco 1 como problema-controle n\~ao-linear. \texttt{make\_moons} com $n=150$, ru\'ido $0{,}15$. O orientador disse na reuni\~ao: "a base da meia-lua voc\^e acha em qualquer lugar a\'i". Era onde esper\'avamos que augmenta\c{c}\~ao R3/R4 ajudasse --- mas j\'a adianto que, no slide da consist\^encia R2/R3/R4, veremos que a m\'etrica diagonal satura em torno de 87\% mesmo ajustada ao r\'otulo verdadeiro, ent\~ao esse ``controle positivo'' vira, na verdade, uma demonstra\c{c}\~ao do limite da diagonal.
\end{frame}

\begin{frame}{Slide 16 --- Bloco 1 --- Oracle Recovery (1/4): tabela do $W$ recuperado}
\footnotesize
Primeira parte do sanity check. Pergunta: se $W$ verdadeiro fosse conhecido, Perceptron e NNLS o recuperariam? Gero dados sint\'eticos com $W$ or\'aculo em tr\^es configura\c{c}\~oes anisotr\'opicas e aprendo $\hat W$ via os dois algoritmos. A tabela mostra resultado m\'edio em 3 seeds. Em x2\_dom (raz\~ao verdadeira 16), NNLS chega a raz\~ao infinita (acerta a dire\c{c}\~ao perfeitamente), cosseno 0,999; Perceptron normaliza para raz\~ao 1,32, cosseno 0,824. Mesmo padr\~ao em x1\_dom e forte\_aniso: NNLS preserva razoes extremas, Perceptron tende a centralizar. Mensagem-chave: ambos atingem fidelidade 100\% --- o sinal classificat\'orio foi capturado --- mas NNLS \'e fiel \`a magnitude exata do $W$ or\'aculo.
\end{frame}

\begin{frame}{Slide 17 --- Bloco 1 --- Oracle Recovery (2/4): visualiza\c{c}\~ao}
\footnotesize
Aqui o gr\'afico grande do painel 2x2: cosseno por configura\c{c}\~ao, fidelidade por configura\c{c}\~ao, raz\~ao $w_1/w_2$ recuperada vs verdadeira, e scatter de $W$ recuperado vs or\'aculo. Visualmente fica claro o padr\~ao da tabela anterior: NNLS pontos pr\'oximos da diagonal (acerta raz\~ao), Perceptron pontos concentrados perto do meio (normaliza). Mensagem para a banca: a discrep\^ancia entre algoritmos n\~ao \'e bug --- \'e diferen\c{c}a metodol\'ogica conhecida. NNLS \'e otim global, Perceptron \'e gradient descent com regulariza\c{c}\~ao impl\'icita pela margem.
\end{frame}

\begin{frame}{Slide 18 --- Bloco 1 --- Oracle Transfer (3/4): tabela de transfer\^encia}
\footnotesize
Segunda parte do sanity check, agora sobre transfer\^encia. Pergunta dual: se aprendo $W$ em uma geometria, ele transfere para outra? Aprendo W em x2\_dom, aplico em x1\_dom, forte\_aniso, Problemas A e E. A tabela mostra cosseno alto entre o $W$ aprendido em origem e o $W$ or\'aculo no destino. Em problemas com geometria similar, transferencia funciona; em Problema A (gaussianas isotr\'opicas), cosseno cai para 0,75 --- a geometria \'e suficientemente diferente. Mensagem: implica\c{c}\~ao direta para Problemas B/C --- sempre recalculamos centr\'oides LOCAIS, n\~ao reaproveitamos do Problema A. A m\'etrica n\~ao \'e independente da geometria do problema.
\end{frame}

\begin{frame}{Slide 19 --- Bloco 1 --- Oracle Transfer (4/4): visualiza\c{c}\~ao}
\footnotesize
Quarto slide do Oracle. Gr\'afico grande mostrando os sub-pain\'eis de transfer\^encia: cada expert origem com suas curvas para os experts destino. Conclus\~ao do sanity check Oracle: ambos algoritmos passam --- recuperam $W$ ($\sim$ fidelidade 100\%) e transferem com cosseno alto entre geometrias similares. Caminho est\'a livre para aplicar o m\'etodo ao LLM nas pr\'oximas se\c{c}\~oes.
\end{frame}

\begin{frame}{Slide 20 --- Bloco 1 --- Problema A descri\c{c}\~ao}
\footnotesize
A Problema A: LLM recebe 150 pontos em zero-shot, responde A ou B. As respostas formam $(X, y_{\text{LLM}})$. Centr\'oides recalculados a partir das respostas do LLM --- ancoramos no que o LLM ACHA serem as classes. Perceptron e NNLS aprendem $W$. Medimos fidelidade.
\end{frame}

\begin{frame}{Slide 21 --- Bloco 1 --- Problema A $W$ e fidelidade}
\footnotesize
Com as 3 seeds, a fidelidade m\'edia da Problema A foi 81{,}3\% com desvio 4{,}4 pontos percentuais. Por seed: 76{,}7, 85{,}3, 82{,}0\%. A acur\'acia do LLM contra o ground truth foi apenas cerca de 33\% --- o LLM inverteu classes em v\'arios casos. Mesmo assim, a m\'etrica reproduz 81\% dessas decis\~oes, demonstrando que H1 mede consist\^encia interna, n\~ao corre\c{c}\~ao. Esse ponto \'e importante para a banca: uma acur\'acia baixa contra a verdade n\~ao invalida a hip\'otese, porque a pergunta \'e se o crit\'erio do LLM \'e est\'avel, n\~ao se ele est\'a certo.
\end{frame}

\begin{frame}{Slide 22 --- Bloco 1 --- Distribui\c{c}\~ao $W$}
\footnotesize
O painel mostra a distribui\c{c}\~ao do $W$ aprendido pelos dois algoritmos nas tr\^es seeds. Boxplot, ratio $w_1/w_2$, scatter de $W$. Estabilidade qualitativa entre seeds.
\end{frame}

\begin{frame}{Slide 23 --- Bloco 1 --- Compara\c{c}\~ao algoritmos}
\footnotesize
Perceptron e NNLS lado a lado: fidelidade, fases B/C, scatter de $W$. Mensagem: os dois algoritmos s\~ao congruentes com cosseno de 0{,}994 nesta execu\c{c}\~ao. Robustez ao m\'etodo de estima\c{c}\~ao.
\end{frame}

\begin{frame}{Slide 24 --- Bloco 1 --- Nossa medida: consist\^encia interna}
\footnotesize
Aqui fa\c{c}o uma decis\~ao metodol\'ogica expl\'icita para a banca. O trabalho mede consist\^encia interna do crit\'erio decis\'orio do LLM: fidelidade na Problema A (quanto $\hat W$ reproduz $y_{\text{LLM}}$), consist\^encia em B/C (quanto $\hat W_A$ prev\^e $y_{\text{LLM}}$ em novos problemas) e acur\'acia LLM-perito na Problema E. Decis\~ao: n\~ao comparamos com ground truth nos problemas sint\'eticos. Raz\~ao: H1 testa se o LLM tem crit\'erio est\'avel, n\~ao se ele est\'a objetivamente correto. Um LLM pode inverter classes e ainda assim ser internamente consistente. Exce\c{c}\~ao: Bloco 3 (peso $\times$ altura) reportar\'a acur\'acia vs r\'otulo real, atendendo solicita\c{c}\~ao expl\'icita do orientador no e-mail das 22:04, ponto 4 --- l\'a h\'a ground truth significativo homem/mulher e classificador \'otimo conhecido.
\end{frame}

\begin{frame}{Slide 25 --- Bloco 1 --- Problema A $W$ por seed e algoritmo}
\footnotesize
Aqui apresento os pesos $W = (w_1, w_2)$ aprendidos para o Problema A, separados por seed e por algoritmo (Perceptron e NNLS). M\'edia das 3 repeti\c{c}\~oes por seed. Os achados-chave:

Perceptron Estruturado: seed 42 $W=(0{,}00023,\ 0{,}00247)$ raz\~ao $0{,}094$; seed 123 $W=(0{,}00030,\ 0{,}00091)$ raz\~ao $0{,}327$; seed 7 $W=(0{,}00005,\ 0{,}00037)$ raz\~ao $0{,}138$.

NNLS: seed 42 $W=(0{,}0357,\ 0{,}1420)$ raz\~ao $0{,}252$; seed 123 $W=(0{,}0466,\ 0{,}1486)$ raz\~ao $0{,}313$; seed 7 $W=(0{,}0394,\ 0{,}1522)$ raz\~ao $0{,}259$.

Antes das observa\c{c}\~oes, vale esclarecer o que \'e \emph{fidelidade na Problema A}: \'e a porcentagem de pontos do Problema A em que o r\'otulo atribu\'ido pela m\'etrica $\hat W$ coincide com o r\'otulo que o LLM atribuiu. \'E medida \emph{in-sample}, ou seja, nos mesmos pontos usados para estimar $\hat W$. Funciona como sanidade: se fosse baixa, o algoritmo n\~ao teria conseguido ajustar uma m\'etrica diagonal aos r\'otulos do LLM, e tudo que viria depois (consist\^encia em B/C) perderia o sentido. Observa\c{c}\~oes importantes para a banca: as magnitudes diferem muito entre algoritmos (Perceptron cerca de cem vezes menor que NNLS), mas ambos apontam na mesma dire\c{c}\~ao --- o cosseno entre os dois $W$ \'e 0{,}994, confirmando que a otimiza\c{c}\~ao inversa \'e invariante \`a escala. Em todas as seeds a raz\~ao $w_1/w_2$ fica abaixo de 1 (cerca de 0{,}25 no NNLS e entre 0{,}09 e 0{,}33 no Perceptron), revelando que o LLM pondera $x_2$ mais que $x_1$, apesar do Problema A ser linearmente separ\'avel em $x_1$ --- pequena anisotropia decis\'oria. Vale registrar que a raz\~ao \'e justamente o par\^ametro mais ru\'idoso entre seeds (coeficiente de varia\c{c}\~ao em torno de 54\%), enquanto a dire\c{c}\~ao permanece est\'avel. Por fim, a fidelidade na Problema A \'e quase id\^entica entre os dois algoritmos, refor\c{c}ando que a estima\c{c}\~ao de $W$ \'e robusta.
\end{frame}

\begin{frame}{Slide 26 --- Bloco 1 --- Problemas B/C tabela consolidada}
\footnotesize
Apresento a tabela das Fases B e C por seed e por algoritmo, m\'edia de 3 repeti\c{c}\~oes por seed, zero-shot. Antes dos n\'umeros, esclare\c{c}o o que \'e consist\^encia: \'e a porcentagem de pontos em que o r\'otulo atribu\'ido pela m\'etrica $\hat W_A$ --- aquela que aprendemos no Problema A --- coincide com o r\'otulo do LLM nos Problemas B e C, que t\^em geometrias rotacionadas. O Kappa de Cohen corrige por concord\^ancia ao acaso. Perceptron por seed: 42 d\'a B 65,0\% / $\kappa$ 0,20 e C 89,0\% / 0,77; 123 d\'a B 66,7\% / 0,28 e C 87,7\% / 0,74; 7 d\'a B 74,3\% / 0,44 e C 85,3\% / 0,69. M\'edia Perceptron: B 68,7\% / 0,31 e C 87,3\% / 0,73. NNLS por seed: 42 d\'a B 68,0\% / 0,27 e C 90,0\% / 0,79; 123 d\'a B 67,0\% / 0,29 e C 89,0\% / 0,77; 7 d\'a B 76,0\% / 0,47 e C 84,0\% / 0,66. M\'edia NNLS: B 70,3\% / 0,35 e C 87,7\% / 0,74. Observa\c{c}\~oes para a banca: H1 \'e sustentada em C (kappa 0,73) mas fica fraca em B (kappa 0,31), ou seja, parcialmente; em todas as seeds o Kappa em C \'e bem maior que em B, indicando que a rota\c{c}\~ao anti-hor\'aria preserva melhor o crit\'erio decis\'orio do LLM. A seed 7 se comporta como outlier sutil --- em B sobe (74\%) mas em C cai mais que as outras. Nota metodol\'ogica: usamos centr\'oides LOCAIS (recalculados em cada problema a partir de $y_{\text{LLM}}$), mais fi\'eis \`a decis\~ao real. O ablativo m\'etrica-contra-Euclidiana em B deu diferen\c{c}a de $+0{,}044$ com intervalo que cont\'em zero (n\~ao significativa): em B, a rota\c{c}\~ao hor\'aria confunde tanto a m\'etrica que ela n\~ao supera a dist\^ancia Euclidiana simples; em C, o ganho \'e claro.
\end{frame}

\begin{frame}{Slide 27 --- Bloco 1 --- Problema B mapa de acertos e erros (seed 42)}
\footnotesize
A pedido do orientador (reuni\~ao 20/05), separei o mapa de acertos e erros em um slide por problema, agora bem maior, para a banca conseguir enxergar ponto a ponto. Este \'e o Problema B, a rota\c{c}\~ao hor\'aria, com a seed 42 como representante --- a tabela anterior j\'a mostrou que ela fica pr\'oxima da m\'edia. Cada ponto verde \'e um acerto da m\'etrica $\hat W_A$ contra o r\'otulo do LLM, cada vermelho \'e um erro. O que importa aqui \'e o padr\~ao geom\'etrico: os erros se concentram na faixa central, exatamente onde as duas gaussianas rotacionadas se sobrep\~oem. A m\'etrica $\hat W_A$ foi aprendida no Problema A, que n\~ao tem rota\c{c}\~ao; a rota\c{c}\~ao hor\'aria de $\pm 1{,}5$ deslocou a fronteira do LLM para uma regi\~ao onde a m\'etrica antiga j\'a n\~ao distingue bem. Por isso a consist\^encia cai para cerca de 68\% e o Kappa para 0,29. Esse mapa \'e a vers\~ao geom\'etrica daquela estat\'istica.
\end{frame}

\begin{frame}{Slide 28 --- Bloco 1 --- Problema C mapa de acertos e erros (seed 42)}
\footnotesize
Este \'e o slide irm\~ao, agora para o Problema C, a rota\c{c}\~ao anti-hor\'aria de mesma magnitude. A compara\c{c}\~ao com o slide anterior \'e o ponto central: aqui os erros aparecem esparsos e isolados, sem concentra\c{c}\~ao em nenhuma regi\~ao do plano. \'E justamente isso que faz a consist\^encia subir para cerca de 90\% e o Kappa para 0,79 --- cerca de tr\^es vezes o Kappa do Problema B, mesmo com magnitude de rota\c{c}\~ao id\^entica. A mensagem para a banca \'e que a dire\c{c}\~ao da rota\c{c}\~ao importa: a anti-hor\'aria preserva muito melhor o crit\'erio decis\'orio que o LLM usou no Problema A do que a hor\'aria. Ver os dois mapas, um em cada slide e em tamanho grande, deixa essa assimetria $\kappa_C \gg \kappa_B$ visualmente \'obvia.
\end{frame}

\begin{frame}{Slide 29 --- Bloco 1 --- Consist\^encia R2/R3/R4: Problema A e meia-lua}
\footnotesize
Tabela por seed comparando R2, R3 e R4 em dois problemas: o A (linear, tr\^es seeds) e a meia-lua (n\~ao-linear). B e C foram omitidos porque s\~ao lineares como A. Peso $\times$ altura fica no Bloco 3. Recapitulando: R2 \'e o espa\c{c}o original $(x_1, x_2)$; R3 acrescenta o produto $x_1 \cdot x_2$; R4 acrescenta $x_1^2$ e $x_2^2$. No Problema A a fidelidade \'e contra os r\'otulos do LLM; o R2 \'e a fidelidade base de duas features (76{,}7 / 85{,}3 / 82{,}0 por seed, m\'edia 81{,}3\%). No Perceptron, R3 fica em 80{,}0 / 83{,}3 / 81{,}3 (m\'edia 81{,}6) e R4 em 78{,}0 / 83{,}3 / 89{,}3 (m\'edia 83{,}6). No NNLS, R3 d\'a 81{,}3 / 86{,}7 / 80{,}7 (m\'edia 82{,}9) e R4 d\'a 85{,}3 / 88{,}0 / 90{,}0 (m\'edia 87{,}8). Leitura para A linear: os ganhos R3/R4 ficam dentro do ru\'ido entre seeds --- augmenta\c{c}\~ao pouco ajuda quando o problema j\'a \'e linearmente separ\'avel. Agora o ponto que mudou nesta execu\c{c}\~ao e que preciso ser honesto com a banca: na meia-lua o LLM colapsou todos os pontos em uma \'unica classe no zero-shot nas TR\^ES seeds --- o log registra ``classe \'unica no LLM, pulando''. Sem dois grupos, n\~ao h\'a r\'otulo do LLM para medir fidelidade, e o pipeline corretamente abortou. Ent\~ao, em vez da fidelidade-contra-o-LLM, reporto o teste-or\'aculo: ajusto a m\'etrica diagonal ao r\'otulo VERDADEIRO da meia-lua. E o resultado \'e revelador: mesmo com o r\'otulo certo, a m\'etrica satura em torno de 86,7\%, cerca de 20 erros em 150, e R3 e R4 n\~ao ajudam --- o NNLS inclusive zera os termos aumentados porque exige pesos n\~ao-negativos. A mensagem \'e que a m\'etrica diagonal tem um teto estrutural na meia-lua: capturar a curvatura em S exigiria um termo cruzado off-diagonal, que a nossa formula\c{c}\~ao diagonal n\~ao tem. Dois achados de uma vez: o zero-shot do gpt-4o-mini na meia-lua \'e inst\'avel, e a m\'etrica diagonal \'e limitada nesse tipo de fronteira --- o que motiva justamente o estudo de caso real do Bloco 3 e entra nas limita\c{c}\~oes.
\end{frame}

\begin{frame}{Slide 30 --- Bloco 1 --- Vi\'es de ordem das classes: tabela}
\footnotesize
Aqui apresento a tabela do experimento de invers\~ao de ordem das classes. Para cada par sem\^antico --- letras, n\'umeros, semantical positivo/negativo e cores --- rodei o mesmo prompt em duas dire\c{c}\~oes: ordem direta e ordem invertida (n\_shot=10, m\'edia de 3 seeds vezes 3 repeti\c{c}\~oes). Pe\c{c}o aten\c{c}\~ao para a coluna ``Fid.\ A'': em A/B caiu de $81{,}3\%$ para $68{,}9\%$ quando invertemos para B/A; em 0/1 subiu de $59{,}8\%$ para $76{,}3\%$ ao virar 1/0; em positivo/negativo a invers\~ao manteve fidelidade alta ($98{,}9\%$ contra $95{,}3\%$); em cores a varia\c{c}\~ao foi pequena ($88{,}7\%$ contra $83{,}8\%$). Leitura para a banca: existe sim um vi\'es de ordem, e ele varia entre $\sim 5$ e $\sim 16$ pontos percentuais na Fid.\ A dependendo do tipo sem\^antico. Letras e n\'umeros s\~ao mais sens\'iveis que cores e r\'otulos sem\^anticos --- o LLM parece ter um v\'inculo mais forte entre certas letras/n\'umeros e ``classe positiva'' que persiste mesmo quando trocamos a ordem do prompt. Por\'em, o vi\'es n\~ao \'e dominante: o critério decis\'orio geral se mant\'em (consist\^encias B e C continuam altas). Portanto o efeito existe, est\'a documentado, mas n\~ao compromete os achados principais.
\end{frame}

\begin{frame}{Slide 31 --- Bloco 1 --- Vi\'es de ordem das classes: gr\'afico}
\footnotesize
Complemento visual da tabela anterior. O gr\'afico \texttt{bloco1\_12\_class\_order\_bias.png} mostra, para cada par sem\^antico, as barras lado a lado da consist\^encia/Fidelidade em ordem direta versus invertida. Visualmente, a banca pode notar que os pares ``A/B vs B/A'' e ``0/1 vs 1/0'' s\~ao os de maior amplitude --- corroborando a leitura num\'erica de que letras e n\'umeros disparam vi\'es maior que cores e r\'otulos sem\^anticos. Para Pos/Neg e Azul/Vermelho as barras s\~ao quase id\^enticas. Conclus\~ao: o vi\'es de ordem est\'a presente, principalmente em pares letra/n\'umero, mas n\~ao explica os achados principais do trabalho.
\end{frame}

\begin{frame}{Slide 32 --- Bloco 1 --- Variantes de prompt: exemplos}
\footnotesize
Antes de mostrar os n\'umeros, exibo o conte\'udo de cada variante para o ponto exemplo $x_1=1{,}23$, $x_2=-0{,}45$ com classes A/B. O \emph{default} \'e o template baseline, sem floreios, s\'o pede a classe. O \emph{geometric} acrescenta a frase explicando que os pontos vivem em um espa\c{c}o m\'etrico euclidiano e que a classifica\c{c}\~ao deve se basear em proximidade espacial aos centros dos clusters. O \emph{CoT} pede que o modelo pense passo a passo e d\^e a resposta final na \'ultima linha. O \emph{tabular} formata as coordenadas em tabela markdown ao inv\'es de prosa. Essas s\~ao perturba\c{c}\~oes pequenas mas controladas; o pr\'oximo slide mostra que essas perturba\c{c}\~oes mudam de fato a m\'etrica recuperada --- ent\~ao o prompt n\~ao \'e neutro.
\end{frame}

\begin{frame}{Slide 33 --- Bloco 1 --- Variantes de prompt: tabela}
\footnotesize
Aqui apresento a tabela com os resultados das 4 variantes de prompt rodadas nas Problemas A--C completas em todas as seeds. Em zero-shot, m\'edia de 3 seeds vezes 3 repeti\c{c}\~oes, configura\c{c}\~ao A/B com features x1/x2. A variante \emph{default} \'e o template baseline. \emph{Geometric} acrescenta contexto espacial expl\'icito (eixos e regi\~oes). \emph{CoT} \'e chain-of-thought, pede racioc\'inio passo a passo. \emph{Tabular} apresenta as coordenadas em tabela markdown. Os n\'umeros: default fid A $81{,}3\%$, cons B $68{,}7\%$ / $\kappa_B$ $0{,}31$, cons C $87{,}3\%$ / $\kappa_C$ $0{,}73$. Geometric fica em $83{,}8$ / $68{,}7$ / $0{,}16$ / $77{,}2$ / $0{,}39$. CoT em $76{,}7$ / $61{,}7$ / $0{,}16$ / $79{,}3$ / $0{,}59$. Tabular \'e o pior: $62{,}0$ / $61{,}4$ / $0{,}13$ / $72{,}2$ / $0{,}21$. Leitura para a banca: o crit\'erio para comparar aqui \'e a \emph{consist\^encia} (kappa em B e C), n\~ao a Fid.\ A isolada. O geometric at\'e empata a Fid.\ A com o default, mas derruba o kappa\_C de 0,73 para 0,39 --- ou seja, prompts mais elaborados n\~ao melhoram a coer\^encia decis\'oria, e o tabular piora tudo. O default puro \'e o que preserva melhor a consist\^encia. CoT, especificamente, n\~ao ajuda em problemas geom\'etricos simples como os nossos --- talvez porque introduza ru\'ido verbal sem informa\c{c}\~ao adicional. A tabela markdown degrada ainda mais, provavelmente porque desloca o LLM para um modo de leitura tabular onde os pares (x1, x2) perdem a intui\c{c}\~ao espacial. Sensibilidade ao prompt \'e portanto \emph{real e n\~ao desprez\'ivel} (queda de at\'e $20$ p.p. em Fid.\ A) --- adotamos o default como refer\^encia e os outros como controle ablativo.
\end{frame}

\begin{frame}{Slide 34 --- Bloco 1 --- Variantes de prompt: gr\'afico}
\footnotesize
Complemento visual da tabela anterior. A imagem \texttt{bloco1\_13\_prompt\_variants.png} mostra, para cada variante de prompt, o scatter dos pesos $W$ aprendidos e as m\'etricas em A, B e C lado a lado. Visualmente fica claro que os pontos de \emph{default} se separam dos pontos de \emph{geometric}, \emph{CoT} e \emph{tabular}. Mensagem para a banca: a configura\c{c}\~ao do prompt n\~ao \'e neutra --- escolhas que parecem inocentes (passar coordenadas em prosa versus tabela versus pedir CoT) modificam mensuravelmente o $W$ que conseguimos recuperar. Por isso a apresenta\c{c}\~ao toda foi padronizada no \emph{default}, e este experimento serve como controle ablativo documentado.
\end{frame}

\begin{frame}{Slide 35 --- Bloco 1 --- Nomes sem\^anticos: tabela}
\footnotesize
Apresento a tabela com o efeito dos nomes sobre o $W$ recuperado, dividida em duas partes. (a) Classes em n\_shot=10 com features $x_1$/$x_2$: A/B d\'a Fid.\ A $81{,}3\%$ e cons B/C $94{,}1/94{,}6\%$; 0/1 colapsa Fid.\ A para $59{,}8\%$; Positivo/Negativo infla Fid.\ A para $98{,}9\%$ mas degrada $\kappa_B$ para $0{,}45$; Azul/Vermelho fica pr\'oximo de A/B. Sem\^anticas n\~ao s\~ao neutras --- o LLM parece ancorar no significado da palavra. (b) Features em zero-shot, classes A/B --- e aqui aviso a banca que esta parte \'e apenas seed 42, portanto ilustrativa: $x_1$/$x_2$ (baseline) Fid.\ A $76{,}7\%$, $\kappa_B$ $0{,}20$, $\kappa_C$ $0{,}77$. \emph{altura/peso} reduz a consist\^encia em B --- $\kappa_B$ cai de $0{,}20$ para $0{,}15$ ---, mas nesta execu\c{c}\~ao, diferente de antes, n\~ao a inverte, n\~ao fica negativo; a Fid.\ A at\'e sobe ($84\%$) e a dire\c{c}\~ao de $W$ \'e preservada. \emph{feature\_1/feature\_2} fica em meio termo ($76\%$, $0{,}32$, $0{,}57$). Mensagem para a banca: os nomes sem\^anticos deslocam a magnitude de $W$ e mexem na consist\^encia, mas o efeito \'e mais brando do que em execu\c{c}\~oes anteriores; por conservadorismo adotamos $x_1$/$x_2$ como baseline em todo o Bloco 1.
\end{frame}

\begin{frame}{Slide 36 --- Bloco 1 --- Nomes sem\^anticos: gr\'aficos}
\footnotesize
Complemento visual: dois pain\'eis lado a lado. \`A esquerda \texttt{bloco1\_14a\_class\_names\_effect.png} com o efeito de mudar nome de classe. \`A direita \texttt{bloco1\_14b\_feature\_names\_effect.png} com o efeito de mudar nome de feature. Visualmente fica claro que o impacto em features (b) \'e bem mais pronunciado que em classes (a) --- argumento concreto para usar nomes neutros ($x_1$/$x_2$, A/B) em todo o Bloco 1.
\end{frame}

\begin{frame}{Slide 37 --- Bloco 1 --- Prompt few-shot (default)}
\footnotesize
Para n\_shot maior que zero, o prompt inclui exemplos rotulados antes do ponto-alvo. Mostro o template literal para n\_shot=5: o LLM v\^e cinco exemplos com $x_1$, $x_2$ e r\'otulo, e depois o ponto novo a classificar. As variantes geometric, CoT e tabular j\'a foram cobertas no slide 31, portanto aqui s\'o exibo o caso default com cinco exemplos para a banca ver como o few-shot \'e constru\'ido na pr\'atica.
\end{frame}

\begin{frame}{Slide 38 --- Bloco 1 --- Resumo Bloco 1}
\footnotesize
Fechando o Bloco 1. Sou honesto com a banca no veredito: H1 fica \emph{parcialmente} sustentada. O Kappa m\'edio em C \'e 0{,}73, acima do limiar de 0{,}5 que fixamos, ent\~ao a consist\^encia se transfere para a rota\c{c}\~ao anti-hor\'aria; mas em B \'e s\'o 0{,}31, abaixo do limiar, pela rota\c{c}\~ao hor\'aria agressiva. Em todas as seeds C supera B, mas n\~ao afirmo H1 como plenamente sustentada. Fidelidade da Problema A em torno de 81\%. Acur\'acia do LLM contra real apenas 33\% --- LLM inverte classes mas continua consistente, o que \'e o cerne da H1. Os dois algoritmos concordam (cosseno 0,994). Sobre R3/R4: na meia-lua a m\'etrica diagonal satura em cerca de 87\% mesmo ajustada ao r\'otulo verdadeiro, ent\~ao augmenta\c{c}\~ao s\'o compensa com n\~ao-linearidade real e dados suficientes --- assunto do Bloco 3. Vieses marginais.
\end{frame}

\begin{frame}{Slide 39 --- Bloco 2 --- Transi\c{c}\~ao Bloco 2}
\footnotesize
Bloco 2: invers\~ao completa de pap\'eis. Um perito externo rotula, o LLM tenta reproduzir o crit\'erio a partir de exemplos few-shot. \'E a Problema E. Problemas: E (linear com $W$ expert) e F (meia-lua com perito = ground truth). Hip\'oteses: H3, H4, H5.
\end{frame}

\begin{frame}{Slide 40 --- Bloco 2 --- Vis\~ao geral Bloco 2}
\footnotesize
Mostro o problema E linear com fronteira do expert quase horizontal, porque $w_2$ \'e cinco vezes maior que $w_1$. A imagem mostra fronteira de decis\~ao, ground truth e margem.
\end{frame}

\begin{frame}{Slide 41 --- Bloco 2 --- Problema E}
\footnotesize
Problema E \'e o antigo "Problema D" da Problema D antiga, agora renomeado. Linear, perito anisotr\'opico $W=[0{,}3, 1{,}5]$. Fronteira quase horizontal. LLM precisa capturar o padr\~ao a partir de exemplos.
\end{frame}

\begin{frame}{Slide 42 --- Bloco 2 --- Problema F meia-lua}
\footnotesize
F usaria a meia-lua com ground truth como perito, para testar se o LLM aprende a fronteira em S s\'o com exemplos --- seria o teste mais duro de H3. Mas j\'a aviso a banca aqui, para n\~ao gerar confus\~ao adiante: nesta execu\c{c}\~ao esse experimento n\~ao p\^ode ser coletado, porque o LLM colapsou em classe \'unica no zero-shot da meia-lua nas tr\^es seeds. Nos pr\'oximos slides eu explico o que fizemos no lugar, o teste-or\'aculo, e por que o teste de H3 em regime n\~ao-linear migra para a base real do Bloco 3.
\end{frame}

\begin{frame}{Slide 43 --- Bloco 2 --- Problema E descri\c{c}\~ao}
\footnotesize
Passo a passo: perito rotula os 150 pontos; divido em treino e teste; para cada $n_{\text{shot}}$ de 0 a 40, seleciono exemplos via estrat\'egia; envio exemplos + ponto de teste em prompt few-shot; me\c{c}o acur\'acia, Kappa, F1 contra perito.
\end{frame}

\begin{frame}{Slide 44 --- Bloco 2 --- Problema E: tabela de resultados}
\footnotesize
Apresento a tabela consolidada dos resultados da Problema E em que o LLM aprende com o perito linear $W=[0{,}3,\ 1{,}5]$ (x2 dominante). M\'edia de 3 seeds vezes 3 repeti\c{c}\~oes. Acur\'acia e Kappa do LLM contra o perito, separados por estrat\'egia (easy, hard, mixed, random) e por n\_shot (0 a 40). Em zero-shot, todas as estrat\'egias degeneram para o mesmo n\'umero ($60{,}9\%$ / $0{,}23$) --- sem exemplos n\~ao h\'a o que escolher. Em n\_shot=5, easy j\'a salta para $84{,}9\%$ enquanto hard fica em $62{,}6\%$. Em n\_shot=20 todos sobem, mas easy ainda lidera com $89{,}8\%$; note que o random \'e o melhor overall entre 5 e 20 exemplos. A invers\~ao acontece em n\_shot=40: hard salta para $97{,}8\%$ e ultrapassa easy ($88{,}3\%$). Mensagem para a banca: H2 (few-shot amplifica) est\'a fortemente sustentada --- saltos de quase $40$ p.p.\ s\'o por adicionar exemplos. E a H4 (exemplos dif\'iceis ajudam mais) sai como \emph{refuta\c{c}\~ao nuan\c{c}ada}: refutada em regime de poucos exemplos, mas confirmada em n\_shot=40. Estrat\'egia n\~ao \'e monotonicamente boa, depende do regime de quantidade de exemplos --- esse \'e um achado novo que vale destacar.
\end{frame}

\begin{frame}{Slide 45 --- Bloco 2 --- Estrat\'egias}
\footnotesize
Quatro estrat\'egias: easy alta margem, hard baixa margem, mixed 50/50, random baseline. A intui\c{c}\~ao de H4 era que hard ensina mais. Vamos ver que \'e mais complexo que isso.
\end{frame}

\begin{frame}{Slide 46 --- Bloco 2 --- Curvas de aprendizado}
\footnotesize
Curvas de acur\'acia versus n\_shot para cada estrat\'egia, m\'edia de 3 seeds vezes 3 repeti\c{c}\~oes. O LLM melhora em todas as estrat\'egias --- H3 sustentada.
\end{frame}

\begin{frame}{Slide 47 --- Bloco 2 --- Por que easy vence hard em poucos exemplos}
\footnotesize
Em modelos com gradiente, hard \'e informativo. LLM in-context n\~ao tem gradiente; usa exemplos como \^ancoras. Easy s\~ao pivots claros: "\'e isto que A parece". Hard s\~ao confusos. Com muitos exemplos, h\'a contexto suficiente para hard ajustar fino o crit\'erio. Achado novo desta execu\c{c}\~ao: din\^amica muda com n\_shot.
\end{frame}

\begin{frame}{Slide 48 --- Bloco 2 --- M\'ultiplos experts}
\footnotesize
Testei tr\^es experts: aniso\_x2 com $W=[0{,}3, 1{,}5]$ fronteira horizontal, aniso\_x1 oposto, e euclidean. As acur\'acias m\'edias: aniso\_x2 vai de 60,9\% em zero-shot a 94,3\% em 40-shot; aniso\_x1 come\c{c}a quase no acaso (48,5\%, o LLM tem vi\'es contra fronteira vertical sem exemplos) e chega a 89,8\%; euclidean vai de 55,2\% a 91,8\%. Todos ultrapassam 90\% com 40 exemplos, e o aniso\_x2 horizontal atinge o melhor patamar final. Mensagem: o LLM aprende qualquer geometria de perito via few-shot --- H3 amplamente sustentada.
\end{frame}

\begin{frame}{Slide 49 --- Bloco 2 --- Dilui\c{c}\~ao}
\footnotesize
Dilui\c{c}\~ao testa se easy dilui o sinal de hard fixos. Fixei 4 hard e adicionei easy progressivamente. Curva monot\^onica crescente --- easy ajuda sempre, n\~ao h\'a dilui\c{c}\~ao. Refor\c{c}a o achado.
\end{frame}

\begin{frame}{Slide 50 --- Bloco 2 --- Vi\'es de ordem dos exemplos few-shot: tabela}
\footnotesize
Aqui mostro a tabela do experimento de vi\'es de ordem. Fixei a estrat\'egia \emph{mixed} e os mesmos exemplos selecionados; variei apenas a ordem em que aparecem no prompt. Testo quatro ordena\c{c}\~oes: class0\_first (todos da classe 0 primeiro), class1\_first (todos da classe 1 primeiro), shuffled (ordem aleat\'oria) e alternating (alternando 0, 1, 0, 1). Roda em n\_shot = 5, 10 e 20. Em n\_shot=5: amplitude m\'axima de $4{,}2$ p.p.\ entre ordena\c{c}\~oes. Em n\_shot=10: amplitude sobe um pouco para $7{,}5$ p.p.\ (shuffled lidera com $87{,}4\%$, class0\_first \'e o pior com $79{,}9\%$). Em n\_shot=20: amplitude colapsa para $2{,}3$ p.p.\ --- todos pr\'oximos de $86\%$. O teste de Kruskal-Wallis n\~ao acusa diferen\c{c}a significativa em nenhum n\_shot. Leitura para a banca: existe sim algum efeito de ordena\c{c}\~ao em n\_shot intermedi\'ario, mas \'e marginal (menor que 8 p.p.\ no pior caso) e some quando se aumenta o n\'umero de exemplos. Recency bias n\~ao explica os achados principais; o crit\'erio decis\'orio do LLM \'e robusto \`a ordem.
\end{frame}

\begin{frame}{Slide 51 --- Bloco 2 --- Vi\'es de ordem dos exemplos few-shot: gr\'afico}
\footnotesize
Complemento visual da tabela anterior. A imagem \texttt{bloco2\_07\_example\_order.png} mostra, para cada n\_shot, as barras das quatro ordena\c{c}\~oes lado a lado. Visualmente fica claro que as barras s\~ao parecidas dentro de cada n\_shot --- corrobora a leitura num\'erica de que recency bias \'e marginal.
\end{frame}

\begin{frame}{Slide 52 --- Bloco 2 --- Problema F (meia-lua): tabela}
\footnotesize
Aqui preciso ser transparente com a banca: nesta execu\c{c}\~ao o experimento de LLM-aprendiz na meia-lua n\~ao p\^ode ser coletado. O LLM colapsou em classe \'unica no zero-shot nas tr\^es seeds --- respondeu tudo como uma classe s\'o ---, e sem duas classes o pipeline few-shot foi abortado. \'E um achado honesto sobre a fragilidade do zero-shot em geometria n\~ao-linear. No lugar, mostro o teste-or\'aculo: ajusto a m\'etrica diagonal ao r\'otulo verdadeiro da meia-lua. Ela satura em cerca de 86,7\%, 20 erros em 150, e R3/R4 n\~ao ajudam. A m\'etrica diagonal tem teto estrutural nesse problema --- precisaria de termo off-diagonal para a curvatura em S.
\end{frame}

\begin{frame}{Slide 53 --- Bloco 2 --- Problema F (meia-lua): colapso visualizado (SVM RBF)}
\footnotesize
Este slide atende o item 7 da reuni\~ao de 20/05: o professor pediu para sobrepor a superf\'icie de separa\c{c}\~ao de um SVM gaussiano \`a fronteira do LLM na meia-lua. A figura tem dois pain\'eis com o mesmo SVM RBF como r\'egua. No painel da esquerda, o SVM treinado no ground truth desenha a fronteira n\~ao-linear de refer\^encia --- \'e o que um classificador competente enxerga nesse problema. No painel da direita, os r\'otulos zero-shot do LLM: todos os 105 pontos de treino receberam a mesma classe, nas tr\^es seeds. N\~ao h\'a fronteira alguma a comparar --- e essa aus\^encia \'e exatamente o resultado. Da\'i tiro duas li\c{c}\~oes: sobre o LLM, sem exemplos o gpt-4o-mini n\~ao percebe as duas semi-luas, uma fragilidade real do zero-shot em geometria n\~ao-linear, n\~ao um bug do pipeline; sobre a m\'etrica, o teste-or\'aculo do slide anterior mostra teto estrutural de 87\% da diagonal. A consequ\^encia pr\'atica \'e que o teste de ``LLM aprendiz'' em regime n\~ao-linear genu\'ino migra para o Bloco 3, na base real peso $\times$ altura, onde o LLM de fato recebe exemplos e aprende, e ali a H3 \'e avaliada com ground truth humano.
\end{frame}

\begin{frame}{Slide 54 --- Bloco 2 --- Baselines cl\'assicos}
\footnotesize
Comparo o LLM com k-NN, Logistic Regression e SVM, todos treinados nos MESMOS exemplos few-shot e avaliados no mesmo conjunto de teste, com o perito linear. E aqui vem um achado honesto e importante: no perito linear os baselines cl\'assicos superam o LLM em TODOS os n\_shot. A Regress\~ao Log\'istica, que casa exatamente com um perito linear, lidera: 84,6\% em 5-shot, subindo a 98,3\% em 40-shot, contra o LLM que vai de 74,7\% a 92,0\%. SVM e k-NN tamb\'em ficam \`a frente. A leitura para a banca \'e direta: quando o crit\'erio \'e trivial e linear, o LLM n\~ao faz melhor que um classificador cl\'assico simples --- a vantagem do LLM, se houver, tem que aparecer em outro regime, e \'e o que investigamos no Bloco 3.
\end{frame}

\begin{frame}{Slide 55 --- Bloco 2 --- LLM vs Perceptron baseline}
\footnotesize
Pedido espec\'ifico do orientador. Para cada $n_{\text{shot}}$, treino um Perceptron Estruturado nos mesmos exemplos. Como a meia-lua colapsou, essa compara\c{c}\~ao aparece na base real peso $\times$ altura, detalhada no Bloco 3. E o resultado se inverteu em rela\c{c}\~ao a execu\c{c}\~oes anteriores: agora \'e o LLM que supera o Perceptron baseline em todos os n\_shot, de +2,2 at\'e +10,7 pontos percentuais. Em 5-shot, por exemplo, o LLM faz 87,4\% contra 76,7\% do Perceptron. N\~ao existe mais o antigo ``Perceptron 90\% contra LLM 83\%''.
\end{frame}

\begin{frame}{Slide 56 --- Bloco 2 --- Resumo Bloco 2}
\footnotesize
Fechando o Bloco 2. H3 sustentada com ganho de cerca de 40 pontos percentuais de 0 para 40-shot no perito linear. H4 refutada em n\_shot pequeno, com signific\^ancia em 5 e 10-shot, mas revertida em 40-shot, onde o hard alcan\c{c}a 0,98 --- diagn\'ostico mais sofisticado do que o esperado. Ponto honesto: os baselines cl\'assicos superam o LLM em todos os n\_shot no perito linear. E o regime n\~ao-linear, que seria a meia-lua, colapsou nesta execu\c{c}\~ao --- ele \'e retomado na base real do Bloco 3. Recency bias e vi\'es de classes marginais.
\end{frame}

\begin{frame}{Slide 57 --- Bloco 3 --- Transi\c{c}\~ao Bloco 3}
\footnotesize
Bloco 3: estudo de caso real, peso $\times$ altura, base do orientador. \'E onde aplicamos tudo o que vimos em dados reais com 100 amostras e fronteira el\'iptica conhecida. Pergunta: priors sem\^anticos do LLM ajudam ou atrapalham?
\end{frame}

\begin{frame}{Slide 58 --- Bloco 3 --- Problema G (peso x altura): apresenta\c{c}\~ao}
\footnotesize
Base enviada pelo orientador no e-mail das 19:15 do dia 30 de abril. 100 amostras, 50 homens e 50 mulheres. Features normalizadas em $[-1, 1]$. Classificador \'otimo bayesiano fornecido pelo professor tem forma el\'iptica.
\end{frame}

\begin{frame}{Slide 59 --- Bloco 3 --- Variantes de prompt na base real}
\footnotesize
Duas variantes. NEUTRA: features = $(x_1, x_2)$, LLM n\~ao sabe que \'e peso e altura. SEM\^ANTICA: features = (peso, altura), LLM ativa priors. Pergunta em ambas: homem ou mulher.
\end{frame}

\begin{frame}{Slide 60 --- Bloco 3 --- Problema G (peso x altura): R2/R3/R4}
\footnotesize
Aprendo $W$ em tr\^es configura\c{c}\~oes. O gr\'afico tem quatro pain\'eis: fidelidade e acur\'acia, pelos dois algoritmos. E, atendendo ao item 3 da reuni\~ao de 20/05, agora mostro tamb\'em o pr\'oprio vetor $W$ estimado, por seed e por algoritmo, na configura\c{c}\~ao de 2 features com prompt sem\^antico. O Perceptron aprende um $W$ quase isotr\'opico --- raz\~ao entre os pesos de peso e altura perto de 1 nas tr\^es seeds ---, enquanto o NNLS aprende um $W$ anisotr\'opico dominado pela altura. Os dois chegam a acur\'acias reais pr\'oximas, o que ilustra um ponto importante para a banca: o $W$ \'e identific\'avel apenas a menos de escala, ent\~ao comparamos dire\c{c}\~oes e raz\~oes, n\~ao valores absolutos.
\end{frame}

\begin{frame}{Slide 61 --- Bloco 3 --- Acur\'acia REAL Problema G (peso x altura)}
\footnotesize
Tabela com prompt sem\^antico. Fidelidade da m\'etrica NNLS sobe de 65\% em R2 para 75\% em R4, enquanto a acur\'acia REAL fica entre 75 e 85\% --- e cai justamente em R4, o sinal do overfitting. Nota importante: essa queda em R4 \'e s\'o do NNLS, porque a restri\c{c}\~ao de pesos n\~ao-negativos zera os termos aumentados; o Perceptron n\~ao cai, fica em torno de 91\%. Em prompt NEUTRO x1/x2, a acur\'acia REAL despenca para 45-49\%, n\'ivel do acaso --- o LLM sem prior sem\^antico est\'a perdido.
\end{frame}

\begin{frame}{Slide 62 --- Bloco 3 --- Paradoxo overfitting}
\footnotesize
Em prompt NEUTRO x1/x2, a fidelidade do NNLS sobe de cerca de 59\% em R2 para 78\% em R4, mas a acur\'acia REAL fica no acaso --- o NNLS vai de 44 para 49\%, e o Perceptron at\'e cai, de 54 para 45\%. A m\'etrica com 4 pesos para 70 amostras de treino decora o ru\'ido das decis\~oes erradas do LLM. \'E o overfitting cl\'assico, e o paradoxo: quanto mais a m\'etrica ``imita'' o LLM, pior ela acerta a verdade. Subir n\_feat s\'o vale com n\~ao-linearidade real E dados suficientes.
\end{frame}

\begin{frame}{Slide 63 --- Bloco 3 --- Problema G (peso x altura): LLM como aprendiz}
\footnotesize
Tabela: LLM com prompt sem\^antico atinge cerca de 96\% de acur\'acia\_metric em 40-shot --- converge para o pr\'oprio crit\'erio ---, e 87\% de acur\'acia real, cerca de 4 erros em 30. E o ponto que mudou: o LLM supera o Perceptron baseline em TODOS os n\_shot, de +2,2 at\'e +10,7 pontos percentuais; em 5-shot faz 87,4\% contra 76,7\%. Diferente de execu\c{c}\~oes anteriores, em que o Perceptron vencia em 5-shot. Com dados rotulados, na base real, o LLM \'e competitivo ou melhor.
\end{frame}

\begin{frame}{Slide 64 --- Bloco 3 --- Problema G: nomes de classe A/B}
\footnotesize
Este slide atende o pedido do e-mail do orientador de 20/05: repetir a classifica\c{c}\~ao pedindo ``A ou B'' em vez de ``homem ou mulher''. A ideia \'e isolar o prior sem\^antico dos nomes de classe, complementando o teste de features neutras. O resultado \'e um dos mais interessantes do Bloco 3: com features sem\^anticas peso e altura, a acur\'acia contra o r\'otulo real despenca de 91,4\% para 12,9\% --- 61 erros em 70. E esse n\'umero, muito abaixo dos 50\% do acaso, n\~ao \'e ru\'ido: \'e invers\~ao sistem\'atica. O LLM continua separando os dois grupos --- se invert\^essemos os r\'otulos, ter\'iamos cerca de 87\% --- mas, sem a \^ancora sem\^antica de ``homem'' e ``mulher'', ele n\~ao tem como saber qual grupo \'e ``A'' e qual \'e ``B'', e a conven\c{c}\~ao saiu oposta \`a do ground truth. Ou seja: o prior sem\^antico n\~ao atua s\'o nas features, atua tamb\'em nos nomes das classes --- \'e ele que ancora o mapeamento entre r\'otulo e grupo. Na Fase E, com exemplos, o LLM adota fielmente a conven\c{c}\~ao dos exemplos que recebe: a concord\^ancia com a m\'etrica-alvo sobe de 74\% para 96\% em 40-shot, e o baseline Perceptron treinado nos mesmos exemplos se comporta da mesma forma --- consist\^encia entre aprendizes, n\~ao um defeito do LLM.
\end{frame}

\begin{frame}{Slide 65 --- Bloco 3 --- Curvas Problema E por variante}
\footnotesize
Curvas separadas por variante, com dois pain\'eis: LLM contra a m\'etrica e LLM contra o real. O prompt sem\^antico domina o neutro em acur\'acia real: fica em torno de 87\% enquanto o neutro patina perto de 47\%. \'E a evid\^encia mais direta de que o prior sem\^antico peso/altura \'e o que carrega o desempenho na base real.
\end{frame}

\begin{frame}{Slide 66 --- Bloco 3 --- LLM vs Perceptron Problema G (peso x altura)}
\footnotesize
Pedido espec\'ifico do orientador. Mostro o gr\'afico com bar chart comparando LLM e Perceptron baseline em cada $n_{\text{shot}}$.
\end{frame}

\begin{frame}{Slide 67 --- Bloco 3 --- Quando o LLM ganha}
\footnotesize
Tabela final, em 10-shot. Na base real o LLM supera o Perceptron baseline nos dois prompts: peso/altura sem\^antico 86,3\% contra 80,7\%, e x1/x2 neutro 47,0\% contra 43,7\%. A diferen\c{c}a \'e que, no neutro, a acur\'acia absoluta \'e baixa para os dois --- perto do acaso ---, enquanto s\'o com o prior sem\^antico ela sobe a 88\%. A meia-lua n\~ao entra na compara\c{c}\~ao porque colapsou. Mensagem honesta para a banca: a vantagem do LLM vem dos priors do treinamento massivo, n\~ao do mecanismo de ICL em si.
\end{frame}

\begin{frame}{Slide 68 --- Bloco 3 --- Visualiza\c{c}\~ao ponto-a-ponto}
\footnotesize
Atende adendo do e-mail das 22:06. Mostro lado a lado peso $\times$ altura com prompt sem\^antico e neutro, ambos 5-shot da seed 42. Os mesmos pontos s\~ao classificados de forma muito diferente. Evid\^encia visual do efeito do prior sem\^antico.
\end{frame}

\begin{frame}{Slide 69 --- Bloco 3 --- Resumo Bloco 3}
\footnotesize
Priors sem\^anticos s\~ao reais: acur\'acia real peso/altura em torno de 88\%, contra cerca de 47\% no x1/x2 neutro. R4, a elipse, recupera a forma te\'orica mas overfit-a com 100 amostras. E nesta execu\c{c}\~ao, na base real, o LLM supera o Perceptron baseline em todos os n\_shot --- resultado revertido em rela\c{c}\~ao ao passado. A meia-lua n\~ao p\^ode ser comparada porque colapsou. A vantagem do LLM vem dos priors, n\~ao do ICL.
\end{frame}

\begin{frame}{Slide 70 --- S\'intese cruzada}
\footnotesize
Tabela final cruzando os 3 blocos. A/B/C lineares: ganho praticamente zero na fidelidade, apenas ru\'ido, como esperado. Meia-lua, pelo teste-or\'aculo: satura em 86,7\% e R3/R4 n\~ao ajudam --- a m\'etrica diagonal tem teto. Peso $\times$ altura no prompt neutro: a fidelidade sobe cerca de 19 pontos percentuais mas a acur\'acia real cai cerca de 16 --- overfit pesado com 100 amostras. O padr\~ao geral: augmenta\c{c}\~ao R3/R4 s\'o compensa com n\~ao-linearidade real e dados suficientes.
\end{frame}

\begin{frame}{Slide 71 --- Robustez entre seeds}
\footnotesize
Tabela: fidelidade Problema A $81{,}3\% \pm 4{,}4$, Consist B $68{,}7\% \pm 5{,}0$, Consist C $87{,}3\% \pm 1{,}9$, kappa B $0{,}31$, kappa C $0{,}73$. Desvios menores ou iguais a 5 p.p.\ entre seeds. A dire\c{c}\~ao de $W$ \'e est\'avel entre seeds, cosseno acima de 0,99, e a raz\~ao $w_1/w_2$ \'e o par\^ametro mais ru\'idoso. H5 sustentada.
\end{frame}

\begin{frame}{Slide 72 --- Testes estat\'isticos principais}
\footnotesize
Zero para 10-shot no Problema B, na consist\^encia (isso \'e H2 no Bloco 1, n\~ao a Fase E): delta $+0{,}255$, intervalo de 0,15 a 0,33, Cohen d de 6,4, efeito enorme. Easy contra Hard em 5-shot: delta $+0{,}223$, intervalo de 0,10 a 0,34 --- \'e a refuta\c{c}\~ao de H4. Em 40-shot inverte para $-0{,}095$, o hard passa o easy. R2 contra R4 na meia-lua, pelo or\'aculo: praticamente zero, satura em 86,7\%. Cosseno entre os dois W: 0,994. Diagn\'ostico de gamma no Perceptron: a figura final\_10 mostra converg\^encia monot\^onica de gamma\_lo.
\end{frame}

\begin{frame}{Slide 73 --- Testes complementares}
\footnotesize
Tr\^es testes adicionais. Point-biserial entre margem e acerto, por seed: $r$ entre 0,28 e 0,39, $p<0{,}001$ --- maior margem, mais acerto. Mann-Whitney comparando margem em acertos contra erros: $p<0{,}001$ nas tr\^es seeds, erros concentrados em baixa margem. Kruskal-Wallis sobre ordem de exemplos: $p \geq 0{,}14$, n\~ao significativo em nenhum n\_shot.
\end{frame}

\begin{frame}{Slide 74 --- S\'intese das hip\'oteses}
\footnotesize
S\'intese por bloco. Bloco 1: H1 \emph{parcialmente} sustentada --- confirmada em C, com kappa 0,73, mas fraca em B, com kappa 0,31 abaixo do limiar de 0,5; e H2 sustentada. Bloco 2: H3 sustentada, H4 refutada em n\_shot pequeno com revers\~ao em 40-shot. Transversal: H5 sustentada, desvio at\'e 5 pontos percentuais. Quatro das cinco sustentadas, com nuances importantes que fa\c{c}o quest\~ao de expor honestamente.
\end{frame}

\begin{frame}{Slide 75 --- Conclus\~oes principais}
\footnotesize
Seis conclus\~oes. Um: crit\'erio decis\'orio est\'avel no LLM, ortogonal a corre\c{c}\~ao. Dois: few-shot amplifica. Tr\^es: o LLM aprende crit\'erio via ICL, mas os baselines cl\'assicos o superam no perito linear em todos os n\_shot. Quatro: easy supera hard em poucos exemplos mas hard reverte em 40-shot --- achado nuan\c{c}ado. Cinco: R3/R4 s\'o ajuda em n\~ao-linear E com dados; na meia-lua a m\'etrica diagonal satura em 87\% mesmo contra a verdade. Seis: a vantagem do LLM em peso $\times$ altura vem dos priors sem\^anticos, e nessa base o LLM supera o Perceptron baseline.
\end{frame}

\begin{frame}{Slide 76 --- Limita\c{c}\~oes}
\footnotesize
M\'etrica diagonal n\~ao captura intera\c{c}\~oes. 2D apenas. Peso $\times$ altura tem s\'o 100 amostras. 1 LLM testado. Sem controle de vers\~ao do LLM entre execu\c{c}\~oes: a coleta usou o alias gpt-4o-mini, sem snapshot datado, entre 30 de junho e primeiro de julho de 2026 --- o provedor pode atualizar o modelo silenciosamente. E h\'a um limite f\'isico da medi\c{c}\~ao que fazemos quest\~ao de declarar: a auditoria mediu 4,753\% de respostas divergentes em consultas id\^enticas a temperatura zero, ent\~ao nenhuma concord\^ancia com o LLM pode passar de uns 95\% --- as nossas m\'etricas s\~ao lidas contra esse teto. Honestidade nas limita\c{c}\~oes \'e o que separa pesquisa de marketing.
\end{frame}

\begin{frame}{Slide 77 --- Agradecimentos}
\footnotesize
Encerro agradecendo o orientador Prof.\ Raul Fonseca Neto, o PPGCC da UFJF e minha fam\'ilia. Estou \`a disposi\c{c}\~ao para perguntas. Tenho dez slides de ap\^endice t\'ecnico cobrindo formula\c{c}\~ao matem\'atica, overfitting, literatura e reprodutibilidade --- inclusive o diagn\'ostico da busca bin\'aria de gamma que o professor pediu.
\end{frame}

\begin{frame}{Slide 78 --- A1 Perceptron deriva\c{c}\~ao}
\footnotesize
Para perguntas matem\'aticas. Formula\c{c}\~ao primal Eq.\ 28: max gamma sujeito \`a restri\c{c}\~ao de margem com relaxa\c{c}\~ao lambda alpha\_i. Regra de corre\c{c}\~ao Eq.\ 29. Atualiza\c{c}\~ao da margem Eq.\ 31. Hiperpar\^ametros do artigo: $\eta = 0{,}001$, $C$ entre 0{,}1 e 1{,}0.
\end{frame}

\begin{frame}{Slide 79 --- A2 NNLS}
\footnotesize
NNLS: $\min \|Aw - b\|^2$ com $w \geq 0$. Linha $i$ de $A$ tem diferen\c{c}as de quadrados aos centr\'oides. $b$ vetor de uns. M\'etodo de conjunto ativo de Lawson-Hanson 1974. Sem hiperpar\^ametros, determin\'istico, \'otimo global.
\end{frame}

\begin{frame}{Slide 80 --- A3 Equival\^encia}
\footnotesize
Perceptron e NNLS s\~ao empiricamente equivalentes nesta execu\c{c}\~ao: cosseno de 0{,}994 entre os $W$'s estimados. Robustez do achado ao m\'etodo de estima\c{c}\~ao.
\end{frame}

\begin{frame}{Slide 81 --- B1 Overfitting}
\footnotesize
Caso emblem\'atico: peso $\times$ altura com prompt neutro $x_1/x_2$, 4 features. Fidelidade do NNLS em torno de 78\% mas acur\'acia REAL apenas cerca de 49\%, no n\'ivel do acaso. A m\'etrica decora o ru\'ido das decis\~oes erradas do LLM. Mitiga\c{c}\~oes: regulariza\c{c}\~ao L1/L2, mais dados, valida\c{c}\~ao cruzada.
\end{frame}

\begin{frame}{Slide 82 --- C1 XAI}
\footnotesize
Conex\~oes com literatura. Ahuja \& Orlin 2001. Schultz \& Joachims 2003. Coelho, Borges, Fonseca Neto CILAMCE 2017. Diferencia\c{c}\~ao: aplica\c{c}\~ao a LLMs, black-box, independente de arquitetura.
\end{frame}

\begin{frame}{Slide 83 --- C2 ICL}
\footnotesize
Brown et al.\ 2020 GPT-3 e few-shot. Min et al.\ 2022 ground truth doesn't matter. Wei et al.\ 2022 emergent abilities. Lyu et al.\ 2023 prompts can\^onicos --- alinhado com refuta\c{c}\~ao de H4 em few-shot pequeno.
\end{frame}

\begin{frame}{Slide 84 --- C3 Metric learning}
\footnotesize
Xing et al.\ 2002 original. Davis et al.\ 2007 ITML. Schultz \& Joachims 2003 SVM-style. Justificativa diagonal Liu et al.\ 2010, Kedem et al.\ 2012.
\end{frame}

\begin{frame}{Slide 85 --- D1 Reprodu\c{c}\~ao + diagn\'ostico gamma}
\footnotesize
Comando \'unico para reprodu\c{c}\~ao. Mostro o gr\'afico do diagn\'ostico da busca bin\'aria de gamma que o professor pediu na reuni\~ao em torno de 9 minutos: gamma\_lo monot\^onico crescente confirma converg\^encia em todas as execu\c{c}\~oes.
\end{frame}

\begin{frame}{Slide 86 --- D2 Limita\c{c}\~oes de reprodutibilidade}
\footnotesize
Honestamente: o LLM evolui silenciosamente entre execu\c{c}\~oes. Temperatura zero N\~AO garante determinismo --- batching em GPU produz diferen\c{c}as residuais --- e n\'os n\~ao s\'o afirmamos isso, n\'os medimos: a auditoria offline das intera\c{c}\~oes, feita pelo script audit\_interactions, encontrou 4,753\% de respostas divergentes entre as 29.116 consultas que se repetiram --- 1.384 flips. Isso quantifica o ru\'ido de medi\c{c}\~ao do protocolo: qualquer fidelidade ou kappa tem um teto pr\'atico de uns 95\%, e \'e assim que interpretamos os n\'umeros. Usamos 3 seeds vezes 3 reps para 9 observa\c{c}\~oes por configura\c{c}\~ao. Estocasticidade residual t\'ipica menor que 2\% entre reps da mesma seed. Para execu\c{c}\~oes futuras, a recomenda\c{c}\~ao \'e fixar o snapshot datado do modelo em vez do alias gpt-4o-mini, que o provedor pode atualizar sem aviso.
\end{frame}

\end{document}
